\section{Conclusion}
\label{sec:conclusion}

Revoking access in the identity system of a multi-tenant RAG application does not remove the copies the application has already made. In a controlled testbed, the index ACL copy, the semantic cache and chat memory each disclosed revoked content on their own, with lifetimes set by the crawl period, the TTL and, for memory, nothing within 48\,h. Response metadata separately allowed near-perfect inference of a hidden document's existence when no defense was in place. Provenance-tainted revocation, which tracks the source set of every derived artifact through summaries and enforces it on every read, together with constant-shape responses, gave zero measured leakage and chance-level inference against the attackers tested, with no change in answer quality and +360\,ms median latency. The evidence is from synthetic data, an exact-match judge, a stub generator, and one small real model. The next steps are a real-data corpus, larger models, a semantic leak judge and adaptive attackers.
